\section{Determinant patterns}

A student who has just learned determinants often tries to compute every determinant by direct expansion. This is understandable, but it is not the best habit. Determinants reward attention to structure. Before expanding, we should ask what kind of matrix we are looking at.

\subsection{Zeros are invitations}

A row or column with many zeros usually tells us where to expand.

\begin{examplebox}
For
\[
A=\begin{pmatrix}
2&0&1&3\\
0&0&5&1\\
4&0&-1&2\\
1&7&0&6
\end{pmatrix},
\]
the second column has only one non-zero entry. Expansion along that column gives
\[
\det A=7\det\begin{pmatrix}
2&1&3\\
0&5&1\\
4&-1&2
\end{pmatrix}.
\]
The matrix itself suggests an efficient first move.
\end{examplebox}

\subsection{Triangular form is a goal}

Triangular matrices have easy determinants, so determinant-preserving row operations can be used to create triangular form.

\begin{examplebox}
Let
\[
A=\begin{pmatrix}
1&1&2\\
2&3&7\\
3&5&12
\end{pmatrix}.
\]
Apply
\[
R_2\leftarrow R_2-2R_1,
\qquad
R_3\leftarrow R_3-3R_1,
\]
and then
\[
R_3\leftarrow R_3-2R_2.
\]
This produces
\[
\begin{pmatrix}
1&1&2\\
0&1&3\\
0&0&0
\end{pmatrix}.
\]
Therefore \(\det A=0\), and the zero row reveals linear dependence.
\end{examplebox}

\subsection{Common factors}

A common factor in one row or column may be taken outside the determinant. For example,
\[
\det\begin{pmatrix}
3x&6x\\
2&5
\end{pmatrix}
=3x\det\begin{pmatrix}
1&2\\
2&5
\end{pmatrix}
=3x.
\]
Such a factor may reveal parameter values for which the determinant vanishes.

\subsection{Vandermonde-type structure}

A famous structured example is the \(3\times3\) Vandermonde determinant
\[
V=\begin{pmatrix}
1&1&1\\
x&y&z\\
x^2&y^2&z^2
\end{pmatrix}.
\]
Its factorization follows from determinant-preserving column operations. Apply
\[
C_2\leftarrow C_2-C_1,
\qquad
C_3\leftarrow C_3-C_1.
\]
Then
\[
\det V
=
\det\begin{pmatrix}
1&0&0\\
x&y-x&z-x\\
x^2&y^2-x^2&z^2-x^2
\end{pmatrix}.
\]
Expanding along the first row gives
\[
\det V
=
\det\begin{pmatrix}
y-x&z-x\\
y^2-x^2&z^2-x^2
\end{pmatrix}.
\]
Using
\[
y^2-x^2=(y-x)(y+x),
\qquad
z^2-x^2=(z-x)(z+x),
\]
we obtain
\[
\det V
=(y-x)(z-x)
\det\begin{pmatrix}
1&1\\
y+x&z+x
\end{pmatrix}.
\]
The remaining determinant is \(z-y\). Hence
\[
\boxed{\det V=(y-x)(z-x)(z-y)}.
\]

The factors describe exactly when two columns become equal. For example, if \(x=y\), then the first two columns coincide and the factor \(y-x\) is zero.

\subsection{Looking before expanding}

\begin{summarybox}
Before computing a determinant, ask:
\begin{itemize}
\item Is the matrix triangular or nearly triangular?
\item Is there a row or column with many zeros?
\item Are two rows or columns equal or proportional?
\item Can determinant-preserving row or column operations create zeros?
\item Is there a common factor in a row or column?
\item If parameters are present, when could dependence occur?
\end{itemize}
\end{summarybox}

\section*{What to remember}

A determinant is not merely a computational ornament. It detects structural information, becomes zero when independence is lost, and changes predictably under row and column operations. Efficient calculations begin by reading the matrix before expanding it.
